\documentclass[conference]{IEEEtran}
\usepackage[T1]{fontenc}
\usepackage{amsmath}
\usepackage{graphicx}
\usepackage{url}
\usepackage{booktabs}
\usepackage{etoolbox}
% Keep IEEE typography while using colons for the abstract and keyword labels.
\patchcmd{\abstract}{---}{:\space}{}{\PackageError{rsdis}{Abstract separator patch failed}{}}
\patchcmd{\IEEEkeywords}{---}{:\space}{}{\PackageError{rsdis}{Keyword separator patch failed}{}}
\title{Energy-Aware Task Allocation for Cooperative Lunar Exploration\\
\large Prototype Design and Evaluation Protocol}
\author{\IEEEauthorblockN{Henrique Correia}
\IEEEauthorblockA{Instituto Superior de\\Engenharia do Porto\\1221309}
\and
\IEEEauthorblockN{Pedro Ribeiro}
\IEEEauthorblockA{Instituto Superior de\\Engenharia do Porto\\1231177}
\and
\IEEEauthorblockN{Francisco Liquito}
\IEEEauthorblockA{Instituto Superior de\\Engenharia do Porto\\1231216}}
\begin{document}
\maketitle
\begin{abstract}
Two rovers explore a lunar test field, share visual marker detections, allocate
inspections, recover unfinished work and return to base. Nearest-rover and
energy-aware policies share battery feasibility constraints. Typed ROS 2 modules,
camera perception and a generated Gazebo world support complete physical
simulation missions. This draft reports functional validation and an evaluation
protocol. Matched policy comparisons remain pending.
\end{abstract}
\begin{IEEEkeywords}
multi-robot systems, lunar exploration, task allocation, ROS 2, Gazebo, energy
\end{IEEEkeywords}

\section{Introduction}
Multi-robot task allocation separates missions, tasks and robot capabilities
\cite{gerkey}. Here, two rovers search separate regions and exchange detections.
An inspection can be assigned to a teammate when travel cost and energy
availability favour that choice.

The research question is whether energy-aware allocation improves completion
and safe return compared with distance-only allocation. The contribution is an
integration and evaluation protocol, without claiming optimal allocation.

\section{Related Software and Scope}
Space ROS supplies a VIPER-inspired rover, Mons Mouton terrain, cameras and
energy plugins \cite{spacedemos,spacesim}, pinned at fixed revisions. The
baseline uses lightweight models. An optional adapter retains upstream rover
geometry and steerable joints with namespaced controllers. Its compact mission
is validated; detailed terrain and physical solar integration remain pending.

ROS 2 Jazzy and Gazebo Harmonic communicate through \texttt{ros\_gz}; OpenCV
performs colour segmentation. Ground-truth odometry isolates allocation from
localisation error. SLAM, rough terrain, mineral identification and physical
solar management are outside the validated scope.

\section{Mission and Architecture}
The mission comprises deployment, sector exploration, visual discovery, allocated
inspection and return to base. Uniquely coloured spheres represent five points
of scientific interest. The rovers have different routes and initial battery
fractions. Success requires every target to be confirmed and both rovers parked
at their respective bases.

\subsection{ROS Components}
Each rover has a custom agent and camera detector. A central coordinator receives
shared detections and status, publishes leased inspection goals and records
telemetry. The agent produces velocity commands, tracks its battery, confirms
nearby targets from local camera observations and handles returning and charging.
A monitoring node publishes RViz markers and paths. A local mission console
shows camera feeds, routes and events, with pause, resume and return controls.

Typed messages separate observations, resource state, ownership and completion.
Sensor and command topics are namespaced; mission topics are shared. All mission
coordinates use the world frame.

\subsection{Visual Perception}
HSV masks and contour filters detect five marker colours. Small, clipped or
strongly non-circular blobs are rejected. Pinhole intrinsics map the two
horizontal silhouette edges to tangent rays. For marker diameter $D$, their
half angular separation $\alpha$ and mean bearing $\beta$, the estimate is
\begin{equation}
z_c \approx \frac{D\cos\beta}{2\sin\alpha}, \qquad
y_c \approx -\frac{D\sin\beta}{2\sin\alpha}.
\end{equation}
The camera offset and closest timestamped rover pose transform the estimate
into world coordinates. This approximation assumes a mostly level field.
Detection error must be measured separately against simulator truth.

An inspection requires the assigned rover to reach its stand-off, face the
marker and receive recent local observations over a dwell period. A teammate's
report alone cannot confirm a task.

\subsection{Control}
Agents follow sector routes and sweep cameras at survey stops. The 80 by 60 m
scene includes mapped rocks, structures and crater exclusions. Bounded A* inflates
hazards, marker estimates and peers by the footprint and a safety margin. A
forward laser detects unlabelled obstacles. Commands are checked for curved
motion and braking; missing range data stops motion. Opposing rovers pass on the
right. Fixed inspection approaches avoid chasing noisy marker estimates. These
safeguards assume the supplied level scene.

\section{Allocation and Energy}
Each target has at most one owner, and each rover receives at most one inspection
at a time. Unique target IDs suppress duplicates. Leases distinguish current
assignments from expired work, and every reassignment receives a new task ID.

Let $d_{ij}$ be rover $i$'s approximate distance to target $j$'s stand-off,
$h_{ij}$ its estimated return distance, $b_i$ its battery fraction and $r$ a
reserve. The synthetic energy estimate is
\begin{equation}
e_{ij}=k_d(d_{ij}+h_{ij})+
k_t\left(\frac{d_{ij}+h_{ij}}{v_{\mathrm{ref}}}+t_s\right).
\end{equation}
Both policies require $b_i \geq e_{ij}+r$. The baseline minimises $d_{ij}$.
The energy-aware policy minimises
\begin{equation}
J_{ij}=d_{ij}+\lambda\frac{e_{ij}+r}{\max(b_i,\epsilon)}.
\end{equation}
The recorded weight $\lambda$ has distance units. This greedy heuristic does not
solve global scheduling and may increase distance.

Energy decreases with measured distance and simulation time. Recharging at a
base is accelerated for practical demonstrations. These fractions are not
calibrated to a real rover. A* detours may exceed the straight-line estimate, so
the agent also checks its own return budget and may release work.

\subsection{Recovery}
A missing heartbeat, unavailable or returning rover, or task timeout releases
ownership. Failed tasks enter a cooldown and have bounded attempts. Results for
an expired task ID are ignored. A service disables and restores a rover for
repeatable failure interventions.

\section{Evaluation Protocol}
Run matched experiments with the same world, initial state, camera settings,
mission duration, resource parameters and failure schedule. Only the policy
changes. Scenarios should include nominal exploration, asymmetric batteries,
temporary rover failure, and the combined energy/failure case.

Record target confirmations, success including return, simulation completion
time, measured distance, synthetic energy consumption, failure-to-reassignment
latency, and target-location error against simulator truth.

Retain failed runs and report success rates alongside completion times. Pair
predefined layouts and random seeds. Save revisions, configuration, telemetry,
events and platform details with each run.

\subsection{Validation Boundary}
Portable tests cover ownership, feasibility, recovery, detours and synthetic
images. Offline rehearsals use shared mission logic, kinematic motion and oracle
detections. Separate Linux gates require real rendered cameras, laser scans,
odometry, assignments and full missions. Their conservative geometric footprint
audit does not measure contact forces or validate physical lunar operation.

\section{Functional Validation and Planned Comparison}
Actual Jazzy/Harmonic runs with software rendering confirm all five targets and
return both rovers (Table~\ref{tab:validation}). The expanded failure run disables
rover 1 at world time 20 s and restores
it at 50 s; its red-marker task transfers to rover 2 and is confirmed by that
rover. Both lightweight runs use source revision \texttt{1317fb3}; the detailed
rover uses \texttt{8bcfb9a}.

\begin{table}[t]
\caption{Functional Gazebo checks (energy policy). Gap: minimum conservative
obstacle clearance; zero overlap samples.}
\label{tab:validation}
\centering
\begin{tabular}{lrr}
\toprule
Model and field & Time (s) & Gap (m)\\
\midrule
Lightweight, compact & 116.0 & 0.5421\\
Lightweight, expanded failure & 656.5 & 0.2827\\
VIPER, compact & 171.5 & 0.5218\\
\bottomrule
\end{tabular}
\end{table}

Both models pass full-circle turning, translation, rendered cameras and laser
checks. Configurations, events and telemetry are retained in the project
repository. These functional checks are not an allocation-policy comparison.

\textbf{The policy comparison remains pending.} Report paired counts, success
rates, the metrics above, platform versions and a recovery timeline. Equal
performance or an energy/time trade-off is a valid finding.

\section{Discussion and Limitations}
Ground-truth localisation, known-size coloured markers, planar control and an
accelerated battery limit external validity. Report sensitivity to the heuristic
weight and error in the return-budget estimate. The central coordinator is a
single point of failure. Further work includes the detailed rover on the expanded
map, uneven terrain, solar charging, terrain-aware navigation and richer perception.

\section{Conclusion}
Gazebo validates compact missions with both rover models and an expanded
lightweight failure mission. The design supports cooperative discovery,
inspection and recovery. Allocation-policy conclusions require matched comparisons.

\begin{thebibliography}{9}
\bibitem{gerkey}
B. P. Gerkey and M. J. Matari\'c,
\emph{A formal analysis and taxonomy of task allocation in multi-robot systems},
International Journal of Robotics Research, vol. 23, no. 9,
pp. 939-954, 2004, doi: 10.1177/0278364904045564.
\bibitem{spacedemos}
Space ROS contributors, \emph{Lunar pole exploration rover demo},
\url{https://github.com/space-ros/demos}, revision
\texttt{4b4eb8447423}, accessed Oct. 29, 2026.
\bibitem{spacesim}
Space ROS contributors, \emph{Simulation assets and plugins},
\url{https://github.com/space-ros/simulation}, revision
\texttt{13226b067d92}, accessed Oct. 29, 2026.
\end{thebibliography}
\end{document}
